
\section{奇点的可去性}\label{sec:4-5}
\zhindex{qidiandekequxing}{奇点的可去性}

\BookRecordIndex{T0214}可去奇点定理断言：若穿孔圆盘上的 \BookRecordIndex{T0003}$J$-全纯曲线取值于紧\BookRecordIndex{T0358}辛流形且具有有限能量，
那么它可以光滑延拓到整个圆盘。当穿孔位于边界上、并且 $u$ 把边界映入紧\BookRecordIndex{T0219}拉格朗日
子流形 $L\subset M$ 时，这一结论同样成立。现在我们给出一个不依赖第~\ref{sec:4-1}
节论证的证明，改用引理~\ref{lem:4-3-1} 的平均值估计。第一步是把第~\ref{sec:4-4}
节的\BookRecordIndex{T0099}等周不等式推广到延拓 $u_\gamma$ 定义在穿孔圆盘上的情形。这与
第~\ref{sec:4-1} 节所用的\BookRecordIndex{T0175}极小曲面单调性有关。我们用 $\ell(\gamma)$ 表示道路或闭路
$\gamma:I\to M$ 的长度。

\begin{lemma}\label{lem:4-5-1}
设 $(M,\omega)$ 是紧\BookRecordIndex{T0358}辛流形，$L\subset M$ 是紧\BookRecordIndex{T0219}拉格朗日子流形，且
$J\in\mathcal J_\tau(M,\omega)$ 是被 $\omega$ \BookRecordIndex{T0372}驯服的\BookRecordIndex{T0088}殆复结构。固定常数 $R>0$，则有：
\begin{enumerate}[label=(\roman*),leftmargin=3em]
\item 设 $u:B_R\setminus\{0\}\to M$ 是穿孔圆盘上的 \BookRecordIndex{T0003}$J$-全纯曲线，并且
$E(u)<\infty$。则存在常数 $r_0\in(0,R]$ 和 $c>0$，使得
\[
  0<r<r_0
  \quad\Longrightarrow\quad
  \frac12\int_{B_r}|du|^2\le c\ell(\gamma_r)^2,
\]
其中 $\gamma_r:\mathbb R/2\pi\mathbb Z\to M$ 是闭路
$\gamma_r(\theta):=u(re^{i\theta})$。

\item 设
\[
  u:(B_R\cap\mathbb H\setminus\{0\},
  B_R\cap\mathbb R\setminus\{0\})\longrightarrow(M,L)
\]
是穿孔半圆盘上的 \BookRecordIndex{T0003}$J$-全纯曲线，并且 $E(u)<\infty$。则存在常数
$r_0\in(0,R]$ 和 $c>0$，使得
\[
  0<r<r_0
  \quad\Longrightarrow\quad
  \frac12\int_{B_r\cap\mathbb H}|du|^2\le c\ell(\gamma_r)^2,
\]
其中 $\gamma_r:[0,\pi]\to M$ 是道路
$\gamma_r(\theta):=u(re^{i\theta})$。
\end{enumerate}
\end{lemma}

\begin{proof}
证明 (i)。选取常数 $c>0$ 和 $\delta>0$，使引理~\ref{lem:4-3-1}(i) 与
定理~\ref{thm:4-4-1}(iii) 的结论对这两个常数成立。再选取
$0<r_0\le R/2$，使得
\[
  4\pi^2cE(u;B_{2r_0})<\delta^2.
\]
对 $0<r\le2r_0$，定义光滑函数
\[
  \varepsilon(r):=E(u;B_r)=\frac12\int_{B_r}|du|^2.
\]
于是 $\varepsilon$ 可以连续延拓到闭区间 $[0,2r_0]$，并且
$\varepsilon(0)=0$。此外，由引理~\ref{lem:4-3-1}(i)，当 $r\le r_0$ 时有
\[
  \frac12|du(re^{i\theta})|^2\le\frac c{r^2}\varepsilon(2r).
\]
现在，$\gamma_r$ 的导数之范数为
\[
  |\dot\gamma_r(\theta)|
  =\frac r{\sqrt2}|du(re^{i\theta})|
  \le\sqrt{c\varepsilon(2r)};
\]
因而当 $r\le r_0$ 时，闭路 $\gamma_r$ 的长度满足
\begin{equation}\label{eq:4-5-1}
  \ell(\gamma_r)
  =\int_0^{2\pi}|\dot\gamma_r(\theta)|\,d\theta
  \le2\pi\sqrt{c\varepsilon(2r)}<\delta.
\end{equation}
所以，当 $r$ 趋于零时，$\ell(\gamma_r)$ 趋于零。

对 $0<\rho<r\le r_0$，令
$u_r:=u_{\gamma_r}:\mathbb D\to M$ 为注~\ref{rem:4-4-3} 中的延拓，并考虑如下球面
$v_{\rho r}:S^2\to M$：先把 $u$ 限制在\BookRecordIndex{T0170}环域 $B_r\setminus B_\rho$ 上，再分别用圆盘
$u_\rho$ 和 $u_r$ 填充边界圆周 $\gamma_\rho$ 与 $\gamma_r$。这个球面可缩，因为它是
一个三维球的边界，而该三维球由 $\rho\le s\le r$ 时的二维圆盘
$u_s:\mathbb D\to M$ 之并构成。因此，对 $0<\rho<r\le r_0$ 有
\[
  E(u;B_r\setminus B_\rho)+\int_{\mathbb D}u_\rho^*\omega
  =\int_{\mathbb D}u_r^*\omega.
\]
令 $\rho\to0$，再使用定理~\ref{thm:4-4-1} 的\BookRecordIndex{T0099}等周不等式，得到
\[
  E(u;B_r)=\int_{\mathbb D}u_r^*\omega\le c\ell(\gamma_r)^2.
\]
这就证明了引理~\ref{lem:4-5-1} 的 (i)。
\end{proof}

\begin{exercise}\label{ex:4-5-2}
证明引理~\ref{lem:4-5-1} 的 (ii)。其证明几乎可以逐字照搬 (i) 的证明。
提示：把短闭路 $\gamma:S^1\to M$ 的\BookRecordIndex{T0202}局部辛作用量，替换为充分短的弧
$\gamma:[0,\pi]\to M$ 的\BookRecordIndex{T0202}局部辛作用量，其中这条弧把两个边界点映入 $L$。
另见 Urs Frauenfelder 的学位论文 \,\srcite{120}。
\end{exercise}

\begin{proof}[定理~\ref{thm:4-1-2} 的证明]
我们证明 (i)，并继续沿用引理~\ref{lem:4-5-1} 证明中的记号。特别地，$c$ 和
$\delta$ 是正数，使引理~\ref{lem:4-3-1}(i) 与定理~\ref{thm:4-4-1}(iii) 的结论
对它们成立。检查定理~\ref{thm:4-4-1} 的证明可知 $c>1/(4\pi)$。此外，
\[
  \begin{aligned}
  \varepsilon(r):=E(u;B_r)
  &=\frac12\int_{B_r}|du|^2\\
  &=\frac12\int_0^r\rho\int_0^{2\pi}
    |du(\rho e^{i\theta})|^2\,d\theta\,d\rho.
  \end{aligned}
\]
由引理~\ref{lem:4-5-1} 的\BookRecordIndex{T0099}等周不等式，当 $r$ 充分小时，
\[
  \begin{aligned}
  \varepsilon(r)
  &\le c\ell(\gamma_r)^2\\
  &=\frac{cr^2}{2}
    \left(\int_0^{2\pi}|du(re^{i\theta})|\,d\theta\right)^2\\
  &\le\pi cr^2\int_0^{2\pi}|du(re^{i\theta})|^2\,d\theta\\
  &=2\pi cr\dot\varepsilon(r).
  \end{aligned}
\]
令 $\mu:=1/(4\pi c)<1$，则上式可以改写为
\[
  \frac{2\mu}{r}\le\frac{\dot\varepsilon(r)}{\varepsilon(r)}.
\]
从 $r$ 到 $r_1$ 积分，得到
\[
  \left(\frac{r_1}{r}\right)^{2\mu}
  \le\frac{\varepsilon(r_1)}{\varepsilon(r)}.
\]
因此
\[
  \varepsilon(r)\le c_1r^{2\mu},
  \qquad c_1:=r_1^{-2\mu}\varepsilon(r_1),
\]
并且由引理~\ref{lem:4-3-1}(i)，当 $\rho$ 充分小时，
\begin{equation}\label{eq:4-5-2}
  |du(\rho e^{i\theta})|^2
  \le\frac c{\rho^2}\varepsilon(2\rho)
  \le\frac{c_2}{\rho^{2-2\mu}},
  \qquad c_2:=cc_1 2^{2\mu}.
\end{equation}
取
\[
  2<p<\frac2{1-\mu},
  \qquad c_3:=2\pi c_2^{p/2},
\]
则当 $r$ 充分小时，
\[
  \begin{aligned}
  \int_{B_r}|du|^p
  &=\int_0^r\int_0^{2\pi}
    \rho|du(\rho e^{i\theta})|^p\,d\theta\,d\rho\\
  &\le c_3\int_0^r\rho^{1-p(1-\mu)}\,d\rho.
  \end{aligned}
\]
由于 $p<2/(1-\mu)$，有
\[
  1-p(1-\mu)>-1,
\]
故这个积分有限。此外，由式~\eqref{eq:4-5-2} 可知，$u$ 在穿孔圆盘上是一致
Hölder 连续的。由于 $M$ 紧（因而完备），$u$ 可以连续延拓到零点。习题
\ref{ex:4-5-4} 又表明，$u$ 在 $\mathbb D$ 上存在弱一阶导数，而且它们与
$\mathbb D\setminus\{0\}$ 上的强导数一致。因此，$u$ 属于从 $\mathbb D$ 到 $M$ 的
索伯列夫映射空间 $W^{1,p}$。在 $u(0)$ 附近取一个局部坐标卡，对 $u$ 应用定理 B.4.1，
即得 $u$ 光滑。
这证明了 (i)；(ii) 的证明见习题~\ref{ex:4-5-5}。定理~\ref{thm:4-1-2} 得证。
\end{proof}

\begin{exercise}\label{ex:4-5-3}
由式~\eqref{eq:4-5-2} 证明 $u$ 是指数为 $\mu$ 的 Hölder 连续映射。
\end{exercise}

\begin{exercise}\label{ex:4-5-4}
设 $u:\mathbb D\to\mathbb R$ 是连续函数，并且它在穿孔圆盘
$\mathbb D\setminus\{0\}$ 上的限制光滑。
再假设 $\partial_su$ 和 $\partial_tu$ 的 $L^p$-范数有界。证明
$u\in W^{1,p}(\mathbb D)$。

提示：证明 $L^p$ 函数 $\partial_su:\mathbb D\to\mathbb R$ 与 $u$ 的\BookRecordIndex{T0293}弱导数一致，即对
每个 $\phi\in C_0^\infty(\mathbb D)$ 都有
\[
  \int_{\mathbb D}\bigl(u\partial_s\phi+(\partial_su)\phi\bigr)=0.
\]
用 Stokes 定理把 $\{|z|\ge\delta\}$ 上的积分表示为 $\{|z|=\delta\}$ 上的积分，
并证明后一个积分在 $\delta\to0$ 时趋于零。
\end{exercise}

\begin{exercise}\label{ex:4-5-5}
证明定理~\ref{thm:4-1-2} 的 (ii)。其证明几乎可以逐字照搬 (i) 的证明。
提示：定义
\[
  \varepsilon(r):=E(u;B_r\cap\mathbb H),
\]
并使用引理~\ref{lem:4-5-1}(ii)（而不是 (i)）。另见 \srcite{302,120}。
\end{exercise}

\begin{exercise}\label{ex:4-5-6}
考察在零点有奇点的全纯映射 $f:\mathbb D\setminus\{0\}\to\mathbb C$。对哪些 $p$，
这样的映射可能属于 $L^p(\mathbb D)$ 或 $W^{1,p}(\mathbb D)$？对取值于 $S^2$ 的全纯
映射考察同一个问题。
把后者看成全纯映射 $\mathbb D\setminus\{0\}\to\mathbb C$ 与嵌入
$\mathbb C\hookrightarrow\mathbb C\cup\{\infty\}$ 的复合，其中
$\mathbb C\cup\{\infty\}$ 是黎曼球面。
\end{exercise}
